% GENERATED FILE. Edit canonical lessons, then run npm run generate:books.
% canonical-source-hash: fnv1a64:37f3e5649139f83c
% canonical-lessons: KA-C192-oralu, KA-C192-onake, KA-C192-miksi, KA-C192-jaradi, KA-C192-turimane

\chapter{Grinding stone, Pestle, Mixer-grinder, Sieve, Coconut grater}
\label{ch:ka-grinding-stone-pestle-mixer-grinder-sieve-coconut-grater}
\hlchaptermodality{192}

\begin{chapteropening}
\textbf{By the end of this chapter:} \emph{I can say a grinding stone, a pestle, a mixer-grinder, a sieve and a coconut grater.}

Complete the last lesson of chapter 192: I can say a grinding stone, a pestle, a mixer-grinder, a sieve and a coconut grater.
\end{chapteropening}

\section[\texorpdfstring{\kn{ಒರಳು} (oraḷu)}{ಒರಳು (oraḷu)}]{\kn{ಒರಳು} (oraḷu) — a grinding stone, a mortar}

\label{lesson:KA-C192-oralu}



\begin{quote}
Before the new one: say the Kannada for a team, then the Kannada for cricket.
\end{quote}

\subsection*{You'll want to know: \kn{ಒರಳು}}
\textbf{\kn{ಒರಳು}} — \emph{oraḷu} — \textquotedblleft{}a grinding stone, a mortar\textquotedblright{}.

\begin{grammarlens}[title={What you've built}]
One more word for this chapter.
\end{grammarlens}

\subsection*{Your turn}
\begin{itemize}
\raggedright
  \item \emph{Say it:} \emph{oraḷu}
  \item \emph{Say it:} \emph{oraḷu}, once more
  \item \emph{Say it:} say \emph{oraḷu}
  \item \emph{Recall:} say the Kannada for a team, then the Kannada for cricket, then say \emph{oraḷu} again
\end{itemize}

\begin{tcolorbox}[breakable,colback=teal!4,colframe=teal!35!black,title={Before you move on}]
What does \kn{ಒರಳು} mean? (\textquotedblleft{}A grinding stone, a mortar\textquotedblright{}.) Say it once more.
\end{tcolorbox}
\section[\texorpdfstring{\kn{ಒನಕೆ} (onake)}{ಒನಕೆ (onake)}]{\kn{ಒನಕೆ} (onake) — a pestle}

\label{lesson:KA-C192-onake}



\begin{quote}
Before the new one: say the Kannada for cricket, then the Kannada for a grinding stone.
\end{quote}

\subsection*{You'll want to know: \kn{ಒನಕೆ}}
\textbf{\kn{ಒನಕೆ}} — \emph{onake} — \textquotedblleft{}a pestle\textquotedblright{}.

\begin{grammarlens}[title={What you've built}]
One more word for this chapter.
\end{grammarlens}

\subsection*{Your turn}
\begin{itemize}
\raggedright
  \item \emph{Say it:} \emph{onake}
  \item \emph{Say it:} \emph{onake}, once more
  \item \emph{Say it:} say \emph{onake}
  \item \emph{Recall:} say the Kannada for cricket, then the Kannada for a grinding stone, then say \emph{onake} again
\end{itemize}

\begin{tcolorbox}[breakable,colback=teal!4,colframe=teal!35!black,title={Before you move on}]
What does \kn{ಒನಕೆ} mean? (\textquotedblleft{}A pestle\textquotedblright{}.) Say it once more.
\end{tcolorbox}
\section[\texorpdfstring{\kn{ಮಿಕ್ಸಿ} (miksi)}{ಮಿಕ್ಸಿ (miksi)}]{\kn{ಮಿಕ್ಸಿ} (miksi) — a mixer-grinder}

\label{lesson:KA-C192-miksi}



\begin{quote}
Before the new one: say the Kannada for a grinding stone, then the Kannada for a pestle.
\end{quote}

\subsection*{You'll want to know: \kn{ಮಿಕ್ಸಿ}}
\textbf{\kn{ಮಿಕ್ಸಿ}} — \emph{miksi} — \textquotedblleft{}a mixer-grinder\textquotedblright{}.

\begin{grammarlens}[title={What you've built}]
One more word for this chapter.
\end{grammarlens}

\subsection*{Your turn}
\begin{itemize}
\raggedright
  \item \emph{Say it:} \emph{miksi}
  \item \emph{Say it:} \emph{miksi}, once more
  \item \emph{Say it:} say \emph{miksi}
  \item \emph{Recall:} say the Kannada for a grinding stone, then the Kannada for a pestle, then say \emph{miksi} again
\end{itemize}

\begin{tcolorbox}[breakable,colback=teal!4,colframe=teal!35!black,title={Before you move on}]
What does \kn{ಮಿಕ್ಸಿ} mean? (\textquotedblleft{}A mixer-grinder\textquotedblright{}.) Say it once more.
\end{tcolorbox}
\section[\texorpdfstring{\kn{ಜರಡಿ} (jaraḍi)}{ಜರಡಿ (jaraḍi)}]{\kn{ಜರಡಿ} (jaraḍi) — a sieve}

\label{lesson:KA-C192-jaradi}



\begin{quote}
Before the new one: say the Kannada for a pestle, then the Kannada for a mixer-grinder.
\end{quote}

\subsection*{You'll want to know: \kn{ಜರಡಿ}}
\textbf{\kn{ಜರಡಿ}} — \emph{jaraḍi} — \textquotedblleft{}a sieve\textquotedblright{}.

\begin{grammarlens}[title={What you've built}]
One more word for this chapter.
\end{grammarlens}

\subsection*{Your turn}
\begin{itemize}
\raggedright
  \item \emph{Say it:} \emph{jaraḍi}
  \item \emph{Say it:} \emph{jaraḍi}, once more
  \item \emph{Say it:} say \emph{jaraḍi}
  \item \emph{Recall:} say the Kannada for a pestle, then the Kannada for a mixer-grinder, then say \emph{jaraḍi} again
\end{itemize}

\begin{tcolorbox}[breakable,colback=teal!4,colframe=teal!35!black,title={Before you move on}]
What does \kn{ಜರಡಿ} mean? (\textquotedblleft{}A sieve\textquotedblright{}.) Say it once more.
\end{tcolorbox}
\section[\texorpdfstring{\kn{ತುರಿಮಣೆ} (turimaṇe)}{ತುರಿಮಣೆ (turimaṇe)}]{\kn{ತುರಿಮಣೆ} (turimaṇe) — a coconut grater}

\label{lesson:KA-C192-turimane}



\begin{quote}
Before the new one: say the Kannada for a mixer-grinder, then the Kannada for a sieve.
\end{quote}

\subsection*{You'll want to know: \kn{ತುರಿಮಣೆ}}
\textbf{\kn{ತುರಿಮಣೆ}} — \emph{turimaṇe} — \textquotedblleft{}a coconut grater\textquotedblright{}.

\begin{grammarlens}[title={What you've built}]
One more word for this chapter.
\end{grammarlens}

\subsection*{Your turn}
\begin{itemize}
\raggedright
  \item \emph{Say it:} \emph{turimaṇe}
  \item \emph{Say it:} \emph{turimaṇe}, once more
  \item \emph{Say it:} say \emph{turimaṇe}
  \item \emph{Recall:} say the Kannada for a mixer-grinder, then the Kannada for a sieve, then say \emph{turimaṇe} again
\end{itemize}

\begin{tcolorbox}[breakable,colback=teal!4,colframe=teal!35!black,title={Before you move on}]
What does \kn{ತುರಿಮಣೆ} mean? (\textquotedblleft{}A coconut grater\textquotedblright{}.) Say it once more.
\end{tcolorbox}
